\section{单调时间、deadline 与竞争完成}

\begin{keyidea}
周期中断不是时间 API。只有把硬件实际周期、整数换算余数、单调表示、等待
排序和唯一唤醒结果连接起来，\code{sleep} 与 timed condition 才具有可以证明
的语义。
\end{keyidea}

\topic{8253/8254 为什么仍在 64 位启动路径上}
IBM PC 用 8253 Programmable Interval Timer 提供可编程周期，PC/AT 后续采用
兼容的 8254。传统通道 0 连接 8259A IRQ0，通道 2 还参与扬声器门控。现代
处理器已经拥有 TSC、Local APIC timer，PC 平台也曾引入 HPET，但兼容历史
保留了 PIT 的 I/O 端口与中断路由。QEMU \code{pc} 机器模拟这条链，使早期
内核不依赖 ACPI、PCI 枚举或固件服务就能取得可重复时钟事件。

PIT 输入并不是十进制整频，而是约 1.193182 MHz。它来自 PC 历史上的
14.31818 MHz 参考频率再除以 12。软件向通道 0 写入 16 位整数除数 \(d\)，
模式 2 周期重装后产生：

\[
f_{\mathrm{irq}}=\frac{1193182}{d}
\]

本项目目标约 1000 Hz，实际写入 \code{0x04A9}，也就是 1193。设备只承诺按
这个整数分频；它从未承诺“一次 IRQ 精确等于一毫秒”。若软件把目标频率
当成硬件事实，短测试看似正常，运行时间增长后误差会稳定累积。

\begin{pdfdiagram}[node distance=8mm and 10mm]
  \node[os hardware] (reference) {14.31818 MHz\\平台参考};
  \node[os state,right=of reference] (pitinput) {\(\div 12\)\\PIT 输入};
  \node[os hardware,right=of pitinput] (counter) {8254 channel 0\\\(\div 1193\)};
  \node[os state,right=of counter] (pic) {IRQ0\\PIC vector 32};
  \node[os state,below=of pic] (runtime) {InterruptRuntime\\推进单调时钟};
  \node[os state,left=of runtime] (deadline) {DeadlineQueue\\解析到期 Thread};
  \draw[os arrow] (reference) -- (pitinput);
  \draw[os arrow] (pitinput) -- (counter);
  \draw[os arrow] (counter) -- (pic);
  \draw[os arrow] (pic) -- (runtime);
  \draw[os arrow] (runtime) -- (deadline);
\end{pdfdiagram}
\diagramnote{硬件只提供分频事件；纳秒、deadline 和 Thread 状态都属于软件契约。}

\topic{clock source、clock event 与 wall clock}
时间子系统容易把三个问题混为一谈：

\begin{itemize}
  \item \textbf{clock source} 回答从某个起点经过了多少时间；
  \item \textbf{clock event} 在某个时刻触发 CPU 进入处理程序；
  \item \textbf{wall clock} 回答日期、时区和民用时间。
\end{itemize}

当前 PIT 周期同时为前两层提供输入：IRQ0 是周期 clock event，输入频率和
除数则允许软件累计经过时间。它没有电池、日期、时区或闰秒知识。调度超时
必须使用 monotonic clock；若使用可以被管理员校准或因 RTC 变化而跳跃的
wall clock，向后调整会延长等待，向前调整会让大量 deadline 同时误到期。

宿主捕获器打印的 \code{[QEMU][T+...ms]} 也不是来宾时间。它从 Python
启动捕获时计时，受 TCG 翻译、宿主负载和串口 PIO 影响。来宾输出的
\code{MONOTONIC\_START\_NS} 才来自自写 PIT、PIC、IDT 与 IRQ0 路径。两者
可以一起定位停顿，却不能比较绝对起点。

\topic{整数余数必须成为持久状态}
设 PIT 输入频率为 \(f\)，实际除数为 \(d\)。一次 tick 对应：

\[
\Delta t=\frac{d\times 10^9}{f}\ \mathrm{ns}
\]

结果通常不是整数。若每次只计算整数商并丢弃余数，\(n\) 次后的软件时间是：

\[
t_n=n\left\lfloor\frac{d\times 10^9}{f}\right\rfloor
\]

真实有理数累计与它之间的差随 \(n\) 增长。v1.13 的
\code{MonotonicClock} 保存整纳秒 \(t\) 和分子余数 \(r\)。一次推进 \(k\)
个 tick 时执行：

\[
(q,r')=\operatorname{divmod}(r+k d 10^9,f),\qquad t'=t+q
\]

且始终满足 \(0\le r'<f\)。余数不是一次计算的临时垃圾，而是下一次计算
所需的时钟状态。逐 tick 推进与一次推进相同总 tick 因而得到同一结果。

\begin{table}[htbp]
\centering
\small
\begin{osgridtabularx}{\textwidth}{lX}
字段 & 不变量与用途\\ \hline
\code{input\_frequency\_hz} & 非零，来自 PIT 硬件规格而不是目标频率\\ \hline
\code{divisor} & 实际写入通道 0 的整数除数\\ \hline
\code{nanoseconds} & 已经提交的整纳秒，只增不减\\ \hline
\code{fraction\_numerator} & 小于输入频率，保存尚未成为一纳秒整数的部分\\ \hline
\code{saturated} & 达到表示上界后保持真，后续推进不再回绕\\ \hline
\end{osgridtabularx}
\caption{单调时钟状态不是一个裸 tick counter}
\end{table}

为什么不用浮点数？freestanding 内核不能假定宿主数学运行时存在；浮点
舍入模式又属于每 Thread 的 FXSAVE 现场。IRQ0 中使用 XMM 或 x87 会把时间
实现与用户扩展现场隔离纠缠。固定宽度整数让宿主模型与目标代码逐位一致，
每个乘法和加法也能在提交前检查。

公式也不能直接当作 64 位机器的求值顺序。实现先把 \(k\) 切成满足
\code{chunk * divisor <= UINT64\_MAX} 的有限分块；每块对
\code{remaining\_cycles * 1e9} 先求商和余数，再把旧余数加入已经小于
\(f\) 的余数。这样周期乘积或“乘积加旧余数”的中间值越界时，仍能保留
最终可表示的精确结果，而不会过早饱和。

\topic{饱和把 64 位边界变成全序}
无符号 64 位纳秒可表达约 584 年。教学系统不会依靠这个运行年限，但仍必须
定义边界。若计数从 \code{UINT64\_MAX} 回绕到零，一个旧的远期 deadline
会突然变成“尚未到达很久”，队列排序也失去时间意义。

实现通过分块消除可避免的中间乘积越界，并检查把最终商加入当前纳秒的范围。
只有时间结果本身不能精确表示时，时钟才提交：

\[
t=\mathrm{UINT64\_MAX},\qquad saturated=true
\]

以后保持不变。用户 \code{SleepFor} 把相对时长转换为绝对 deadline 时也
采用饱和加法。这样极大时长表示“最远可表达未来”，不会伪装成一个已经过去
的小数字。

\topic{为什么内核只接受绝对 deadline}
若等待层层保存“还剩多少”，每次虚假唤醒、信号或 Mutex 重获都必须重新
扣除已经经过的时间。某一层忘记扣除会延长总等待，多层整数换算又会增加误差。
v1.13 在接口边界只转换一次：

\[
deadline=\operatorname{saturating\_add}(now,duration)
\]

Kernel、ThreadScheduler 和 DeadlineQueue 都保存同一个绝对纳秒。Condition
被唤醒但谓词仍不成立时，下一轮继续使用原 deadline，而不是重新取得完整
时长。这也为 v1.14 signal 中断等待留下明确语义。

\topic{DeadlineQueue 的身份与稳定顺序}
当前系统最多 512 个 Thread，每个 Thread 同时只能处于一个 Blocked 关系，
因此最多拥有一个活动 deadline。队列为每个 Thread 索引保留直接可定位槽，
活动槽再按如下键组成有序关系：

\[
(absolute\ deadline,\ insertion\ sequence)
\]

sequence 使相同 deadline 的 Thread 有稳定顺序。测试参考模型因此能给出唯一
预期，而不是接受任意排列。当前插入为 \(O(n)\)，查看和移除最早项为
\(O(1)\)。在 512 的明确硬上限内，它比尚未需要的复杂时间轮更容易审计。
以后可替换容器，但不能改变稳定排序、按 Thread 取消和统计语义。

队列故意不保存用户地址、Futex 指针或业务 WaitQueue 指针。ThreadScheduler
已经拥有 Thread 当前等待关系；timer 层只保存 Thread 身份。到期时调度器
从 Thread 读取权威 WaitQueue membership。这样 \code{munmap}、exec 和
ProcessExit 不会在 timer 层留下第二份悬空用户地址索引。

\topic{一个等待同时属于两张图}
timed wait 让 Thread 同时出现在两种关系中：

\begin{pdfdiagram}[node distance=12mm and 18mm]
  \node[os state] (thread) {Thread\\Blocked};
  \node[os state,below left=of thread] (wait) {WaitQueue\\业务条件};
  \node[os state,below right=of thread] (deadline) {DeadlineQueue\\绝对时刻};
  \node[os state,below=25mm of thread] (ready) {Thread\\Ready};
  \draw[os arrow] (thread) -- node[left,font=\scriptsize]{membership} (wait);
  \draw[os arrow] (thread) -- node[right,font=\scriptsize]{deadline} (deadline);
  \draw[os arrow] (wait) --
      node[below left,font=\scriptsize]{condition/close/cancel} (ready);
  \draw[os arrow] (deadline) --
      node[below right,font=\scriptsize]{timeout} (ready);
\end{pdfdiagram}
\diagramnote{两条边不是两个完成权；同一 scheduler lock 只允许一次 Blocked 到 Ready。}

登记也必须是事务。若先加入 WaitQueue、解锁后才登记 deadline，通知可能在
中间获胜，随后代码又给已经 Ready 的 Thread 留下 deadline。反向顺序则可能
让 IRQ0 到期后找不到业务队列。v1.13 的
\code{BlockCurrentThreadUntil} 在 scheduler irq-save lock 下同时验证当前
Running Thread、deadline 与两种容量，再发布 Blocked。

\topic{条件与超时的单赢家竞争}
通知路径和 IRQ0 到期路径都取得同一把 scheduler lock：

\begin{enumerate}
  \item 胜者确认 Thread 仍为 Blocked；
  \item 从业务 WaitQueue 摘除 Thread；
  \item 普通通知取消 deadline，超时则消费已经弹出的 deadline；
  \item 写入唯一 \code{WakeReason}；
  \item 把 Thread 接入一次 Ready queue；
  \item 败者随后观察到状态不再是 Blocked，只能返回未完成。
\end{enumerate}

普通 wake、Timeout、Thread terminate、unmap cancellation、exec 和
ProcessExit 因而共享同一个状态提交点。稳定统计满足：

\[
active=schedules-expirations-cancellations
\]

系统结束时 \(active=0\)，所以
\(schedules=expirations+cancellations\)。只检查“程序退出码为零”无法发现
一个已退出 Thread 仍留在 deadline queue；账本把这种残留变成发布失败。

\topic{IRQ0 到用户 RAX 的完整控制流}
每次 IRQ0 先推进 MonotonicClock，再处理调度 tick 和 PIC EOI。随后
ProcessRuntime 在 scheduler lock 下反复查看最早 deadline，直到队首仍在
未来。每个到期项按等待条件解析：

\begin{itemize}
  \item Sleep 把保存用户 frame 的返回寄存器写为成功；
  \item private futex 把返回寄存器写为 \code{TIMED\_OUT}；
  \item futex WaitQueue 清空后释放对应 key 槽；
  \item 任一 Thread 变为 Ready 时提交 \code{need-resched}。
\end{itemize}

系统调用阻塞后，C++ 函数本身不会普通返回到原用户 frame。该 Thread 日后
被选中时，汇编从它保存的 176 字节 UserContext 恢复寄存器。因此超时路径
必须修改保存 frame 的 RAX，而不是某个 IRQ0 局部变量。

\topic{idle 是 sleep 正确性的硬件证据}
最后一个 Running Thread 调用 Sleep 后可能没有 Ready Thread。内核切回永久
地址空间和启动栈，执行相邻的 \code{sti; hlt; cli}：

\begin{enumerate}
  \item \code{sti} 开放 IRQ；
  \item \code{hlt} 在没有事件时停止取指，不消耗忙等时间片；
  \item PIT IRQ0 唤醒 CPU，推进单调时钟并解析 deadline；
  \item 中断返回后的 \code{cli} 恢复调度临界区；
  \item 安全返回边界选择新 Ready Thread。
\end{enumerate}

若 deadline 解析错误地要求当前 Thread 有效，或只处理来自 CPL3 的 IRQ，
系统进入 idle 后就永远不会醒。64 MiB bootstrap 只有一个用户 Thread，
正好是最直接的非忙等证明；额外放置一个轮询线程只能掩盖错误。

\topic{timed futex 仍要关闭 compare-and-block 窗口}
\code{WaitPrivateFutexUntil} 不能先在锁外比较用户 word，再单独登记 deadline。
正确事务在 wake 使用的同一 scheduler lock 下完成：

\begin{enumerate}
  \item 再次从用户地址复制 32 位 word；
  \item 与 expected 比较，不同则返回 value changed；
  \item 读取当前单调纳秒，deadline 已到则返回 timeout；
  \item 取得 \((AddressSpaceId,user\ VA)\) futex entry；
  \item 同时登记 WaitQueue 与 deadline；
  \item 最后才释放锁。
\end{enumerate}

用户复制若失败、entry 容量不足或调度登记失败，都在解锁前逆序释放本次取得
的空 entry。Timeout、普通 wake、unmap、exec 和 ProcessExit 最终都必须使
空 entry 可回收，否则重复超时会缓慢耗尽 512 个 key。

\topic{ConditionVariable 返回前为什么必须重新加锁}
ConditionVariable 不保存业务谓词。通知只说明“谓词可能变化”，调用者应当：

\begin{lstlisting}
mutex.Lock();
while (!predicate) {
    result = condition.WaitUntil(mutex, deadline);
    if (result == TimedOut) {
        break;
    }
}
\end{lstlisting}

运行库先观察单调 sequence，释放 Mutex，再进入 timed futex。无论通知、
timeout 还是失败，它都先重新取得 Mutex，才返回强类型结果
\code{ConditionSatisfied}、\code{TimedOut} 或 \code{Failed}。否则调用者
会误以为仍在临界区内访问谓词。sequence 与内核锁内二次比较共同消除
unlock 到 block 之间的丢失通知。

\topic{固定 ABI 与边界}
v1.13 在既有 0--53 系统调用后增加：

\begin{table}[htbp]
\centering
\begin{osgridtabularx}{\textwidth}{lclX}
接口 & 编号 & 返回 & 语义\\ \hline
\code{GetMonotonicTime} & 54 & \code{uint64\_t} & 当前单调纳秒\\ \hline
\code{SleepUntil} & 55 & 状态 & 阻塞到绝对 deadline，已到立即成功\\ \hline
\code{WaitPrivateFutexUntil} & 56 & 状态 & compare-and-block 加绝对 deadline\\ \hline
\end{osgridtabularx}
\caption{v1.13 固定宽度用户 ABI}
\end{table}

超时结果冻结为 \code{-51}。接口没有 \code{time\_t}、\code{long} 或平台
相关宽度，也没有 Kernel 指针。wall clock、RTC、日期、时区、tickless 和
高精度 timer 均明确不在本阶段。

\topic{四层测试为什么不能相互替代}
时间错误既可能来自算术，也可能来自队列竞争或真实中断路径，因此证据必须
从可控时钟模型逐层延伸到 QEMU 中的 PIT、PIC、IDT 和用户态系统调用。

\topic{单元层}
\code{MonotonicClock} 测试精确控制输入频率、除数和 tick 数，验证余数携带、
逐步/批量等价、长跨度和饱和。DeadlineQueue 测试稳定同 deadline 顺序、
重复 Thread 拒绝、任意取消、批量到期、容量与结构校验。

\topic{集成层}
虚拟时钟、ThreadScheduler 与两个 WaitQueue 组合，证明 deadline 前不早醒、
Timeout 只发生一次、condition 先赢会取消 deadline，以及败者重复操作没有
副作用。最终 Thread 回收后队列与统计都必须守恒。

\topic{固定种子随机层}
参考模型在 100000 步中随机执行 schedule、cancel、advance 与 expire。每一步
比较活动集合、最早项、稳定顺序和统计；失败输出固定种子与步骤。它覆盖人工
样例难以穷举的长交错，又保持完全可复现。

\topic{QEMU 整机层}
\filepath{/bin/time\_probe} 经真实 PIT、PIC、IDT、idle、系统调用和调度路径
验证 20 ms Sleep 不早醒、10 ms futex timeout、通知先于 200 ms deadline
以及 10 ms condition timeout。64 MiB 因每 Process 只有一个 Thread，明确
输出单线程 profile；256 MiB 与 64 GiB 必须创建 notifier Thread，让普通
通知真实赢过 deadline。

\topic{日志必须观察边界而不改变热路径}
一千赫兹 IRQ 若每次写串口，会把 PIO 延迟注入被测调度时序。v1.13 只在用户
事务完成和系统回收边界打印：

\begin{lstlisting}
[OS][USER][TIME] MONOTONIC_START_NS=0x...
[OS][USER][TIME] MONOTONIC_WAKE_NS=0x...
[OS][USER][TIME] SLEEP_VERIFIED
[OS][USER][TIME] FUTEX_TIMEOUT_VERIFIED
[OS][USER][TIME] CONDITION_WON_BEFORE_DEADLINE
[OS][USER][TIME] CONDITION_TIMEOUT_VERIFIED
[OS][KERNEL][TIME] ACTIVE_DEADLINES=0x0000000000000000
\end{lstlisting}

start/wake 数字证明来宾时间前进，一次性语义标记证明操作已经提交，最终汇总
证明没有残留。三类证据互补，不能用大量逐 tick 文本替代结构不变量。

\topic{调试顺序与下一阶段输入}
Sleep 不醒时先检查 idle IRQ 是否仍解析 deadline；不能先扩大宿主 timeout。
时间漂移时检查余数是否跨 tick 保留；时间变小时检查乘加回绕和饱和。一个
Thread Ready 两次时检查通知与到期是否共用 scheduler lock，普通 wake 是否
取消 deadline。futex 容量缓慢下降时检查 timeout 和取消路径是否释放空 key。

v1.14 将加入 signal，使 condition、deadline 与异步通知成为三方竞争。signal
必须复用本节的 Blocked 状态与唯一 WakeReason，不能建立一条绕开 WaitQueue
和 deadline 的平行唤醒路径。这样时间阶段不是一次性 Sleep 功能，而是后续
全部可中断阻塞语义的基础。
